\section{Introduction}

Loss of spoken language due to neurological injury or disease is among the most disabling consequences of neurological disorders, profoundly limiting social participation and employment. Brain-computer interface (BCI) systems offer a path toward restoring communication by translating cortical activity into speech output. Intracranial recordings obtained from epilepsy surgery patients, and especially subdural electrocorticography (ECoG), provide cortical signals with sufficient spatial and temporal resolution to support speech decoding research, and a growing body of work has shown that ECoG can be used to read out phonemic, articulatory, and spectrotemporal correlates of spoken and attempted speech~\cite{pasley2012reconstructing,bouchard2014control,mugler2014direct,herff2015brain,flinker2015redefining,wang2021dynamics}. Despite these advances, translating neural activity into intelligible, natural-sounding speech at BCI-compatible latencies remains an open problem, and two challenges dominate: the limited amount of paired neural and speech data that can be collected from a single participant, and the high variability of speech production even for a single speaker producing the same word.

Earlier work on ECoG-to-speech decoding relied on linear models that reached modest Pearson correlation coefficients (PCC) of around $0.6$ between original and reconstructed spectrograms while offering interpretability and low data requirements. More recent deep learning approaches have improved on this baseline using convolutional, recurrent, and hybrid architectures~\cite{akbari2019towards,anumanchipalli2019speech,makin2020machine}. These systems differ in both the intermediate latent representation used to model speech and the resulting audio quality. For example, cortical activity can be decoded into an articulatory kinematic space, which is then transformed into speech waveforms with strong accuracy but non-natural voice quality~\cite{anumanchipalli2019speech}. Other approaches generate more naturalistic audio using wavenet-style vocoders~\cite{akbari2019towards}, generative adversarial networks, or unit-selection, but often at the cost of decoding accuracy. In parallel, clinical-grade speech BCIs have achieved remarkable progress: intracortical microelectrode arrays now support high-vocabulary text decoding at conversational speeds~\cite{willett2023highperformance}, and ECoG-based systems have been used to restore text, audio, and avatar speech output in individuals with severe paralysis~\cite{moses2021neuroprosthesis,metzger2023highperformance,silva2024bilingual}. Most of these systems, however, rely on non-causal operations that can exploit auditory feedback unavailable in real-time settings, limiting their applicability to closed-loop prostheses.

A second gap in the literature is the scarcity of publicly available code. Most state-of-the-art pipelines are not released in a form that allows independent replication or head-to-head comparisons of decoder architectures, causality, electrode density, or hemispheric coverage. Almost all prior studies also focus on left-hemisphere language-dominant cortex, even though many patients who would benefit from a speech BCI have left-hemisphere damage and could in principle be served by a right-hemisphere prosthesis~\cite{lorach2023walking}.

To address these issues, we propose a neural speech decoding framework centered on two components: (i) an \emph{ECoG Decoder} that maps cortical activity to a compact set of interpretable acoustic speech parameters (pitch, voicing, formant frequencies, loudness, broadband noise characteristics), and (ii) a novel fully differentiable \emph{Speech Synthesizer} that maps these parameters to spectrograms. A companion speech-to-speech auto-encoder, comprising a Speech Encoder and the same Speech Synthesizer, is pre-trained on speech signals alone to produce reference speech parameters that supervise ECoG Decoder training. Because the Speech Synthesizer is differentiable, spectral reconstruction loss can be back-propagated end-to-end. The low dimensionality of the parameter space and the availability of reference parameters together mitigate the data scarcity problem. We provide three ECoG Decoder architectures -- 3D ResNet, 3D Swin Transformer, and LSTM -- each supporting both causal and non-causal temporal operations, and apply the framework to 48 neurosurgical patients undergoing epilepsy monitoring.

Our contributions can be summarized as follows:
\begin{itemize}
\item A two-stage training pipeline that combines a differentiable Speech Synthesizer with an ECoG Decoder, supervised by reference speech parameters derived from a speech-only auto-encoder.
\item A direct comparison of convolutional, recurrent, and transformer backbones in both causal and non-causal configurations across 48 ECoG participants, showing that 3D ResNet achieves the highest PCC ($0.804$ non-causal, $0.798$ causal) and 3D Swin closely follows ($0.793$ non-causal, $0.796$ causal).
\item Evidence that causal ResNet and causal Swin models match their non-causal counterparts (Wilcoxon signed-rank $p=0.1022$ and $p=0.2794$, respectively), while causal LSTM is significantly inferior ($p=0.0007$).
\item The first systematic evaluation of right-hemisphere speech decoding with modern deep models, with right-hemisphere PCC of $0.790$ (ResNet) and $0.781$ (Swin) and no statistically significant difference from the left hemisphere.
\item Evidence that low-density clinical grids support decoding performance comparable to hybrid high-density grids.
\item An occlusion-based contribution analysis that identifies distinct cortical substrates engaged by causal versus non-causal decoders.
\item An openly released pipeline to facilitate replication and direct comparison across research groups.
\end{itemize}

\section{Related Work}

\paragraph{Linear and classical ECoG-to-speech decoding.}
Early ECoG-to-speech work used linear regression from cortical high-gamma activity onto acoustic spectrograms or phonetic features~\cite{pasley2012reconstructing,mugler2014direct,herff2015brain}. These models achieved roughly $0.6$ PCC between original and reconstructed spectrograms and remained interpretable with modest data requirements, but suffered from limited accuracy and intelligibility. A complementary line of work characterized the cortical representation of speech in auditory and sensorimotor cortex, showing that ventral sensorimotor cortex encodes spoken vowel acoustics~\cite{bouchard2014control}, that Broca's area coordinates cortical networks for articulation~\cite{flinker2015redefining,wang2021dynamics}, and that auditory cortex predictively tracks speech envelopes~\cite{golumbic2013visual,yildiz2016predictive}. These findings motivate decoder architectures that leverage peri-sylvian cortex, including superior temporal, inferior frontal, pre-central and post-central gyri.

\paragraph{Deep-learning ECoG-to-speech decoding.}
More recent work has turned to deep neural networks. Anumanchipalli et al.~\cite{anumanchipalli2019speech} introduced a recurrent pipeline that decodes cortical activity into articulatory kinematic representations and then synthesizes audible speech; the kinematic latent space supports partial cross-participant transfer but yields a non-naturalistic voice. Akbari et al.~\cite{akbari2019towards} used a deep neural network to reconstruct speech from human auditory cortex and improved digit-level intelligibility by 65\% over a linear baseline. Makin et al.~\cite{makin2020machine} applied an encoder-decoder recurrent network to decode ECoG directly to text and reported word error rates as low as 3\% on a closed sentence set. Subsequent clinical systems have extended these ideas to individuals with severe paralysis: Moses et al.~\cite{moses2021neuroprosthesis} reported real-time word and sentence decoding from sensorimotor cortex of a person with anarthria; Willett et al.~\cite{willett2023highperformance} achieved 62 words per minute and a 9.1\% word error rate on a 50-word vocabulary using intracortical microelectrode arrays; and Metzger et al.~\cite{metzger2023highperformance} demonstrated synchronous text, speech audio, and avatar control from high-density ECoG in a paralyzed participant. Silva et al.~\cite{silva2024bilingual} further showed that a shared articulatory representation supports bilingual English--Spanish decoding. These systems establish the clinical feasibility of speech BCIs but have been reported in case-study format with limited code availability, and most of them rely on non-causal operations at decoding time. Our framework is complementary: it evaluates three decoder backbones in causal and non-causal modes, on a cohort of 48 participants, and is released as an open pipeline~\cite{anumanchipalli2019speech,makin2020machine,moses2021neuroprosthesis,metzger2023highperformance,silva2024bilingual}.

\paragraph{Backbone architectures.}
We build on three well-established deep-learning backbones. 3D ResNet~\cite{he2016deep} serves as a strong convolutional baseline and is adapted to treat stacks of ECoG frames as volumetric inputs. The 3D Swin Transformer extends the Swin family~\cite{liu2021swin,liu2022video} to spatio-temporal tokens, using 2$\times$2$\times$2 patches and patch-merging stages that down-sample in both time and space. The LSTM backbone~\cite{hochreiter1997lstm} represents the widely used recurrent decoding family. Training of all decoders uses the Adam optimizer~\cite{kingma2015adam}.

\paragraph{Speech quality evaluation and speech parameter estimation.}
Speech decoding quality is typically evaluated through spectrogram correlation and intelligibility metrics. We use the Pearson correlation coefficient together with the Short-Time Objective Intelligibility (STOI) family of measures, which operates on one-third octave bands and temporal envelopes~\cite{taal2011algorithm}, including the envelope-power-spectrum variant~\cite{relanoiborra2016sepsm}. Reference pitch and formant targets are computed using the Praat analysis toolkit~\cite{boersma2001praat}.

\paragraph{Neural prostheses beyond speech.}
Finally, our work sits within a broader body of evidence that cortical signals can drive clinically useful motor prostheses, from early real-time arm reach decoding~\cite{wessberg2000realtime} to recent brain-spine interfaces that restored natural walking after spinal cord injury~\cite{lorach2023walking}. These results argue that well-engineered decoders and carefully selected cortical signals can generalize across effectors, and motivate pursuing speech prostheses that exploit both low-density clinical grids and bilateral cortical coverage.
